\chapter{CAR Representation and Fock Space}
\label{ch:car_fock}

\section{Introduction and Abstract Setting}

In quantum field theory, the fundamental kinematics of fermions is governed by the Canonical Anticommutation Relations (CAR). Before assigning spatial or momentum coordinates to fermionic modes, the algebraic theory of fermions must be established over a generic, abstract set of single-particle states.

Let $\iota$ be a non-empty, finite index set equipped with a linear order:
\begin{equation}
\iota = \{1, 2, \dots, d\}, \quad d \coloneqq \#\iota \ge 1.
\end{equation}
The choice of a total linear order $\le$ on $\iota$ is physically unobservable but mathematically essential: it fixes the sign conventions acquired during the permutation of fermionic creation operators, eliminating sign ambiguities.

Typical physical realizations of $\iota$ include:
\begin{enumerate}
    \item The discrete centered Brillouin zone $\iota = \Lambda^* = \{-h+1, \dots, h\}$, linearly ordered by the standard integer order $\le$ (Chapters~\ref{ch:lattice_band} and~\ref{ch:fermions}).
    \item Multi-species or spinful chiral fermions $\iota = \{R, L\} \times \{\uparrow, \downarrow\} \times \Lambda^*$, endowed with the lexicographic product order.
\end{enumerate}

\section{The Two-Tiered CAR Hierarchy: Algebraic vs. Hilbert Setting}

To separate pure algebraic identities from metric and inner-product properties, we distinguish between two tiers of CAR representations.

\begin{definition}[Algebraic CAR Representation — Tier 1]
Let $V$ be a complex vector space. An \emph{Algebraic CAR representation} over $\iota$ on $V$ consists of two families of linear operators $c, c^\dagger : \iota \to \mathrm{End}_{\mathbb{C}}(V)$ satisfying, for all $i, j \in \iota$:
\begin{align}
\{ c_i, c_j \} &\coloneqq c_i c_j + c_j c_i = 0, \label{eq:car_ann_ann} \\
\{ c_i^\dagger, c_j^\dagger \} &\coloneqq c_i^\dagger c_j^\dagger + c_j^\dagger c_i^\dagger = 0, \label{eq:car_cre_cre} \\
\{ c_i, c_j^\dagger \} &\coloneqq c_i c_j^\dagger + c_j^\dagger c_i = \delta_{i, j} I_V. \label{eq:car_ann_cre}
\end{align}
\end{definition}

Without assuming any inner product, adjoint operation, or topological structure on $V$, the algebraic relations alone imply fundamental commutation theorems.

\begin{theorem}[Purely Algebraic CAR Identities]
\label{thm:algebraic_car}
Let $(c, c^\dagger)$ be an Algebraic CAR representation on $V$.
\begin{enumerate}
    \item \textbf{Bilinear Commutator Identity:} For any modes $a, b, c, d \in \iota$:
    \begin{equation}
    \left[ c_a^\dagger c_b, c_c^\dagger c_d \right] = \delta_{b, c}\, c_a^\dagger c_d - \delta_{a, d}\, c_c^\dagger c_b.
    \end{equation}
    \item \textbf{Number Operator Commutators:} For $n_i \coloneqq c_i^\dagger c_i$ and any $j \in \iota$:
    \begin{align}
    [n_i, c_j^\dagger] &= \delta_{i, j} c_j^\dagger, \\
    [n_i, c_j] &= -\delta_{i, j} c_j.
    \end{align}
    \item \textbf{Idempotence and Number Commutativity:}
    \begin{equation}
    n_i^2 = n_i, \qquad [n_i, n_j] = 0 \quad \text{for all } i, j \in \iota.
    \end{equation}
\end{enumerate}
\end{theorem}

\begin{proof}
We verify each algebraic identity:
\begin{enumerate}
    \item We expand the commutator $[c_a^\dagger c_b, c_c^\dagger c_d] = c_a^\dagger c_b c_c^\dagger c_d - c_c^\dagger c_d c_a^\dagger c_b$. 
    Using $c_b c_c^\dagger = \delta_{b, c} I - c_c^\dagger c_b$ on the first term:
    \[
    c_a^\dagger c_b c_c^\dagger c_d = c_a^\dagger (\delta_{b, c} I - c_c^\dagger c_b) c_d = \delta_{b, c} c_a^\dagger c_d - c_a^\dagger c_c^\dagger c_b c_d.
    \]
    Using $c_c^\dagger c_a^\dagger = - c_a^\dagger c_c^\dagger$ and $c_d c_b = - c_b c_d$ on the second term:
    \[
    c_c^\dagger c_d c_a^\dagger c_b = c_c^\dagger (\delta_{a, d} I - c_a^\dagger c_d) c_b = \delta_{a, d} c_c^\dagger c_b - c_c^\dagger c_a^\dagger c_d c_b = \delta_{a, d} c_c^\dagger c_b - (-c_a^\dagger c_c^\dagger)(-c_b c_d).
    \]
    Subtracting the two expressions, the quartic products $-c_a^\dagger c_c^\dagger c_b c_d$ cancel identically, leaving:
    \[
    [c_a^\dagger c_b, c_c^\dagger c_d] = \delta_{b, c} c_a^\dagger c_d - \delta_{a, d} c_c^\dagger c_b.
    \]
    
    \item Setting $a = b = i$ and $c = j$ in the bilinear identity gives $[n_i, c_j^\dagger c_d] = \delta_{i, j} c_i^\dagger c_d - \delta_{i, d} c_j^\dagger c_i$. Contracting with a test state or using $\{c_i, c_j^\dagger\} = \delta_{i,j}$:
    \[
    n_i c_j^\dagger = c_i^\dagger c_i c_j^\dagger = c_i^\dagger (\delta_{i, j} I - c_j^\dagger c_i) = \delta_{i, j} c_i^\dagger - c_i^\dagger c_j^\dagger c_i = \delta_{i, j} c_j^\dagger + c_j^\dagger c_i^\dagger c_i = \delta_{i, j} c_j^\dagger + c_j^\dagger n_i,
    \]
    which gives $[n_i, c_j^\dagger] = \delta_{i, j} c_j^\dagger$. The annihilation commutator follows analogously.
    
    \item For idempotence: $n_i^2 = c_i^\dagger c_i c_i^\dagger c_i = c_i^\dagger (I - c_i^\dagger c_i) c_i = c_i^\dagger c_i - (c_i^\dagger)^2 c_i^2$. Since $(c_i^\dagger)^2 = 0$, $n_i^2 = n_i$. Finally, $[n_i, n_j] = [c_i^\dagger c_i, c_j^\dagger c_j] = \delta_{i, j} c_i^\dagger c_j - \delta_{i, j} c_j^\dagger c_i = 0$.
\end{enumerate}
\end{proof}

\begin{definition}[Hilbert Space CAR Representation — Tier 2]
Let $V$ be a finite-dimensional complex Hilbert space equipped with positive-definite inner product $\langle \cdot, \cdot \rangle$. A \emph{Hilbert CAR representation} over $\iota$ on $V$ is an Algebraic CAR representation $(c, c^\dagger)$ that additionally satisfies \emph{Hilbert Adjoint Compatibility}:
\begin{equation}
c_i^\dagger = (c_i)^\dagger \quad \text{for all } i \in \iota,
\end{equation}
where $(A)^\dagger$ denotes the unique Hilbert space adjoint satisfying $\langle u, A v \rangle = \langle A^\dagger u, v \rangle$.
\end{definition}

\begin{remark}[Formalization Note]
In the verified Lean formalization (\texttt{Bosonize.A02}), the separation between \texttt{AlgebraicCAR} and \texttt{CAR} is strictly maintained. Adjoint compatibility is an independent geometric axiom: if $(c, c^\dagger)$ is an algebraic CAR pair and $S$ is an invertible non-unitary operator, the similarity transform $\tilde{c}_i \coloneqq S c_i S^{-1}$ and $\tilde{c}_i^\dagger \coloneqq S c_i^\dagger S^{-1}$ preserves all anticommutators, but $\tilde{c}_i^\dagger \neq (\tilde{c}_i)^\dagger$. Thus, adjointness is an essential condition that fixes the physical Hilbert space metric.
\end{remark}

\section{The Concrete Occupation Fock Space}

We now construct the canonical Fock space representation of the CAR algebra.

\begin{definition}[Occupation Basis and Fock Space Carrier]
Let $\iota$ be a finite set. The fermionic Fock space $\mathcal{F}(\iota)$ is the complex Euclidean space indexed by the power set $\mathcal{P}(\iota)$ of all subsets of $\iota$:
\begin{equation}
\mathcal{F}(\iota) \coloneqq \ell^2(\mathcal{P}(\iota)) \equiv \bigoplus_{S \subseteq \iota} \mathbb{C} |S\rangle.
\end{equation}
For each subset $S \subseteq \iota$, the basis vector $|S\rangle$ is the indicator state representing the configuration where exactly the single-particle modes in $S$ are occupied.
\end{definition}

\begin{proposition}[Dimension and Inner Product]
\begin{enumerate}
    \item The complex dimension of the Fock space is:
    \begin{equation}
    \dim_{\mathbb{C}} \mathcal{F}(\iota) = 2^{\#\iota} = 2^d.
    \end{equation}
    \item Basis states form an orthonormal basis under the standard Euclidean inner product:
    \begin{equation}
    \langle S | T \rangle = \delta_{S, T} \quad \text{for all } S, T \subseteq \iota.
    \end{equation}
    \item The state $|\emptyset\rangle$ is the unique normalized fermionic vacuum state.
\end{enumerate}
\end{proposition}

\section{Exact Fermionic Sign Mechanics}

To define creation and annihilation operators on $\mathcal{F}(\iota)$ that satisfy anticommutation, we must account for the permutation phase acquired when inserting or removing a fermion.

\begin{definition}[Preceding Count and Fermionic Sign]
Let $\iota$ be linearly ordered by $\le$. For any mode $i \in \iota$ and any occupation subset $S \subseteq \iota$:
\begin{enumerate}
    \item The \emph{preceding count} $\sigma(i, S)$ is the number of occupied modes in $S$ strictly preceding $i$:
    \begin{equation}
    \sigma(i, S) \coloneqq \#\{j \in S \mid j < i\}.
    \end{equation}
    \item The \emph{fermionic sign} $\epsilon(i, S)$ is the sign factor:
    \begin{equation}
    \epsilon(i, S) \coloneqq (-1)^{\sigma(i, S)} \in \{+1, -1\}.
    \end{equation}
\end{enumerate}
\end{definition}

In constructive mathematics, to avoid underflow errors when computing counts in $\mathbb{N}$ (where $0 - 1 = 0$), all set operations are formulated additively.

\begin{lemma}[Additive Atomic Counts]
\label{lem:atomic_counts}
Let $i, j \in \iota$ and $S \subseteq \iota$.
\begin{enumerate}
    \item If $j \notin S$, then:
    \begin{equation}
    \sigma(i, S \cup \{j\}) = \sigma(i, S) + \begin{cases} 1, & \text{if } j < i, \\ 0, & \text{if } j \ge i. \end{cases}
    \end{equation}
    \item If $j \in S$, then:
    \begin{equation}
    \sigma(i, S) = \sigma(i, S \setminus \{j\}) + \begin{cases} 1, & \text{if } j < i, \\ 0, & \text{if } j \ge i. \end{cases}
    \end{equation}
    \item Self-invariance: $\sigma(i, S \cup \{i\}) = \sigma(i, S)$ and $\sigma(i, S \setminus \{i\}) = \sigma(i, S)$.
\end{enumerate}
\end{lemma}

\begin{proof}
The preceding set $\{k \in S \cup \{j\} \mid k < i\}$ is the disjoint union of $\{k \in S \mid k < i\}$ and $\{j\}$ if $j < i$, or simply $\{k \in S \mid k < i\}$ if $j \ge i$. Taking set cardinalities yields the additive formula. Part (2) is the inverse statement, and part (3) follows because $i < i$ is false.
\end{proof}

\begin{theorem}[Fermionic Sign Exchange Identities]
\label{thm:sign_exchange}
For any distinct modes $i \neq j$ in $\iota$:
\begin{enumerate}
    \item \textbf{Creation Pair:} If $i, j \notin S$:
    \begin{equation}
    \epsilon(j, S)\, \epsilon(i, S \cup \{j\}) = - \epsilon(i, S)\, \epsilon(j, S \cup \{i\}).
    \end{equation}
    \item \textbf{Annihilation Pair:} If $i, j \in S$:
    \begin{equation}
    \epsilon(j, S)\, \epsilon(i, S \setminus \{j\}) = - \epsilon(i, S)\, \epsilon(j, S \setminus \{i\}).
    \end{equation}
    \item \textbf{Mixed Pair:} If $i \notin S$ and $j \in S$:
    \begin{equation}
    \epsilon(j, S)\, \epsilon(i, S \setminus \{j\}) = - \epsilon(i, S)\, \epsilon(j, S \cup \{i\}).
    \end{equation}
\end{enumerate}
\end{theorem}

\begin{proof}
Because $\iota$ is linearly ordered and $i \neq j$, either $j < i$ or $i < j$.
\begin{enumerate}
    \item Suppose $j < i$. By Lemma~\ref{lem:atomic_counts}(1):
    \[
    \sigma(i, S \cup \{j\}) = \sigma(i, S) + 1 \implies \epsilon(i, S \cup \{j\}) = (-1)^{\sigma(i, S)+1} = - \epsilon(i, S).
    \]
    On the other hand, since $i > j$, inserting $i$ does not affect the count of elements preceding $j$:
    \[
    \sigma(j, S \cup \{i\}) = \sigma(j, S) \implies \epsilon(j, S \cup \{i\}) = \epsilon(j, S).
    \]
    Therefore:
    \[
    \epsilon(j, S) \epsilon(i, S \cup \{j\}) = \epsilon(j, S) (-\epsilon(i, S)) = - \epsilon(i, S) \epsilon(j, S \cup \{i\}).
    \]
    If $i < j$, the exact same argument applies with the roles of $i$ and $j$ interchanged.
    \item Parts (2) and (3) follow by the same parity discrepancy using Lemma~\ref{lem:atomic_counts}(2).
\end{enumerate}
\end{proof}

\section{Creation and Annihilation Operators on Fock Space}

\begin{definition}[Fermionic Operators on Basis States]
For each mode $i \in \iota$, the annihilation operator $c_i$ and creation operator $c_i^\dagger$ are defined on basis states $|S\rangle$ by:
\begin{align}
c_i |S\rangle &\coloneqq \begin{cases}
\epsilon(i, S) |S \setminus \{i\}\rangle, & \text{if } i \in S, \\
0, & \text{if } i \notin S,
\end{cases} \label{eq:fock_c} \\
c_i^\dagger |S\rangle &\coloneqq \begin{cases}
\epsilon(i, S) |S \cup \{i\}\rangle, & \text{if } i \notin S, \\
0, & \text{if } i \in S,
\end{cases} \label{eq:fock_cdag}
\end{align}
and extended linearly to all states in $\mathcal{F}(\iota)$.
\end{definition}

\begin{theorem}[Verification of the Concrete CAR Representation]
\label{thm:fock_car_verification}
The operators $\{c_i, c_i^\dagger \mid i \in \iota\}$ satisfy:
\begin{enumerate}
    \item \textbf{Adjoint Compatibility:} $c_i^\dagger = (c_i)^\dagger$ under the Euclidean inner product.
    \item \textbf{Anticommutation Relations:}
    \begin{equation}
    \{c_i, c_j\} = 0, \qquad \{c_i^\dagger, c_j^\dagger\} = 0, \qquad \{c_i, c_j^\dagger\} = \delta_{i, j} I.
    \end{equation}
\end{enumerate}
In particular, $\mathcal{F}(\iota)$ is a concrete Hilbert CAR representation.
\end{theorem}

\begin{proof}
We verify both properties on arbitrary basis states $|S\rangle, |T\rangle$:
\begin{enumerate}
    \item For adjointness, we compute the matrix elements $\langle S | c_i | T \rangle$ and $\langle c_i^\dagger S | T \rangle$:
    \begin{itemize}
        \item If $i \in T$ and $S = T \setminus \{i\}$, then $c_i |T\rangle = \epsilon(i, T) |S\rangle$, so $\langle S | c_i | T \rangle = \epsilon(i, T)$. Also, $i \notin S$ and $S \cup \{i\} = T$, so $c_i^\dagger |S\rangle = \epsilon(i, S) |T\rangle$. By Lemma~\ref{lem:atomic_counts}(3), $\sigma(i, S) = \sigma(i, T)$, so $\epsilon(i, S) = \epsilon(i, T)$. Thus $\langle c_i^\dagger S | T \rangle = \epsilon(i, S) = \epsilon(i, T) = \langle S | c_i | T \rangle$.
        \item In all other cases, both matrix elements vanish identically.
    \end{itemize}
    Since this holds for all basis states, $c_i^\dagger = (c_i)^\dagger$.

    \item For anticommutators:
    \begin{itemize}
        \item If $i = j$: $c_i^2 |S\rangle = 0$ because if $i \in S$, $c_i |S\rangle \propto |S \setminus \{i\}\rangle$, which no longer contains $i$, so applying $c_i$ again yields $0$. Similarly, $(c_i^\dagger)^2 = 0$. For the mixed term: if $i \in S$, $c_i^\dagger c_i |S\rangle = |S\rangle$ and $c_i c_i^\dagger |S\rangle = 0$; if $i \notin S$, $c_i^\dagger c_i |S\rangle = 0$ and $c_i c_i^\dagger |S\rangle = |S\rangle$. In both cases, $(c_i c_i^\dagger + c_i^\dagger c_i) |S\rangle = |S\rangle = \delta_{i, i} |S\rangle$.
        \item If $i \neq j$: when $i, j \notin S$, $(c_i^\dagger c_j^\dagger + c_j^\dagger c_i^\dagger)|S\rangle = \big( \epsilon(j, S) \epsilon(i, S \cup \{j\}) + \epsilon(i, S) \epsilon(j, S \cup \{i\}) \big) |S \cup \{i, j\}\rangle$. By Theorem~\ref{thm:sign_exchange}(1), the sum of signs is $0$. The same holds for $\{c_i, c_j\} = 0$ and $\{c_i, c_j^\dagger\} = 0$ using parts (2) and (3) of Theorem~\ref{thm:sign_exchange}.
    \end{itemize}
\end{enumerate}
\end{proof}

\section{Observables, Parity, and the Hopping Algebra}

\begin{definition}[Number and Parity Observables]
\begin{enumerate}
    \item \textbf{Local Number Operator:} $n_i \coloneqq c_i^\dagger c_i$. On basis states:
    \begin{equation}
    n_i |S\rangle = \begin{cases} |S\rangle, & \text{if } i \in S, \\ 0, & \text{if } i \notin S. \end{cases}
    \end{equation}
    \item \textbf{Total Particle Number Operator:} $\hat{N}_{\mathrm{tot}} \coloneqq \sum_{i \in \iota} n_i$. On basis states:
    \begin{equation}
    \hat{N}_{\mathrm{tot}} |S\rangle = (\#S)\, |S\rangle.
    \end{equation}
    \item \textbf{Global Fermion Parity Operator $\Gamma$:} Defined as the ordered product over the sorted index set $\iota = \{i_1 < i_2 < \dots < i_d\}$:
    \begin{equation}
    \Gamma \coloneqq \prod_{k=1}^d (I - 2 n_{i_k}).
    \end{equation}
\end{enumerate}
\end{definition}

\begin{proposition}[Properties of Global Fermion Parity]
\label{prop:parity}
The parity operator $\Gamma$ satisfies:
\begin{enumerate}
    \item Action on basis states:
    \begin{equation}
    \Gamma |S\rangle = (-1)^{\#S} |S\rangle.
    \end{equation}
    \item Involution and self-adjointness: $\Gamma^2 = I$ and $\Gamma^\dagger = \Gamma$.
    \item Anticommutation with fermionic generators:
    \begin{equation}
    \{\Gamma, c_i\} = 0, \qquad \{\Gamma, c_i^\dagger\} = 0 \quad \text{for all } i \in \iota.
    \end{equation}
\end{enumerate}
\end{proposition}

\begin{proof}
For any $i \in \iota$, $(I - 2n_i)|S\rangle = -|S\rangle$ if $i \in S$, and $+|S\rangle$ if $i \notin S$. Since all $n_i$ commute with each other, the product evaluates to $\prod_{i \in S} (-1) = (-1)^{\#S}$.
Since $(-1)^{\#S} \in \{+1, -1\}$, $\Gamma^2 |S\rangle = (-1)^{2\#S} |S\rangle = |S\rangle$, and $\Gamma$ is self-adjoint.
Finally, $c_i$ reduces $\#S$ by $1$ if $i \in S$, so $\Gamma c_i |S\rangle = (-1)^{\#S - 1} c_i |S\rangle = - c_i (-1)^{\#S} |S\rangle = - c_i \Gamma |S\rangle$.
\end{proof}

\begin{definition}[Hopping Operator]
For any modes $i, j \in \iota$, the hopping operator $T_{i, j}$ is defined by:
\begin{equation}
T_{i, j} \coloneqq c_i^\dagger c_j.
\end{equation}
\end{definition}

\begin{theorem}[Hopping Lie Algebra]
The hopping operators satisfy:
\begin{enumerate}
    \item Adjoint relation: $T_{i, j}^\dagger = T_{j, i}$.
    \item Commutation relations forming the Lie algebra $\mathfrak{u}(d)$:
    \begin{equation}
    [T_{a, b}, T_{c, d}] = \delta_{b, c} T_{a, d} - \delta_{a, d} T_{c, b}.
    \end{equation}
\end{enumerate}
\end{theorem}

\begin{proof}
This is an immediate consequence of Theorem~\ref{thm:algebraic_car}(1) and adjoint compatibility $c_i^\dagger = (c_i)^\dagger$.
\end{proof}

The hopping operators $T_{i, j}$ form the exact algebraic building blocks for the kinetic Hamiltonian and the density operators in the chapters that follow.
